\chapter{StateGraph: State, Nodes and Edges}
\label{ch:stategraph}

\begin{quote}\itshape
Nodes do not return the state. They return a diff, and something else decides
what that means.
\end{quote}

\noindent
This is the chapter the rest of Part~III depends on. LangGraph is not a
convenience wrapper for calling a model in a loop; it is a state machine executed
under a bulk-synchronous parallel model borrowed from Pregel, and almost every
surprising behaviour --- why a parallel write raises, why a message sent in a
step is not visible in that step, why a checkpoint lands where it does --- falls
out of that model rather than from any particular API.

\chapterstub{
\textbf{Argument.} Reducers are the load-bearing concept, not nodes. A field
without a reducer is overwritten; a field with one is merged; and two parallel
nodes writing to a field without a reducer is an \api{InvalidUpdateError} rather
than a race. Understanding the merge rule is understanding the framework.

\textbf{Sections.} (8.1) When to drop from \api{create\_agent} to a raw
\api{StateGraph}, with the four symptoms that justify it. (8.2) The Pregel/BSP
model and supersteps: what runs together, when updates land, where the checkpoint
boundary is. (8.3) The state schema: \code{TypedDict}, \code{Annotated}, reducers,
\api{add\_messages} and its deduplication by id, \api{RemoveMessage}. (8.4)
Writing a custom reducer, and the properties one must have. (8.5) Input and
output schemas; private channels for state that must not leak to the caller.
(8.6) Nodes and edges; \code{START}, \code{END}, conditional edges. (8.7)
\api{Command}: state update and jump in one return, and \code{Command.PARENT} for
crossing a subgraph boundary. (8.8) \api{Send}: dynamic fan-out and map-reduce
over a list. (8.9) Subgraphs: state mapping, checkpointer scope, an agent as a
node. (8.10) Runtime and \code{context\_schema} --- the three-way distinction
between runtime context, graph state and config, which is a common point of
confusion.

\textbf{Listings.} A three-node graph with a conditional edge and a retry loop;
a custom reducer with its test; \api{Send}-based map-reduce over N documents; the
parallel-write error, reproduced and then fixed; a subgraph with explicit state
mapping.

\textbf{Diagrams.} \code{lg-superstep} --- three supersteps drawn as swimlanes,
showing buffered delivery. \code{lg-reducer-merge} --- two parallel updates
merging, with and without a reducer. \code{lg-send-fanout} --- map-reduce over a
document list. \code{lg-state-vs-context-vs-config} --- the three scopes as a
Venn-style diagram.

\textbf{Mini-project.} Stage 05: the agent rebuilt as an explicit graph, with
\api{Send} fan-out across retrieved documents.

\textbf{Source topics covered.} Part~IV 4.1--4.8.
}
